\documentclass[11pt]{article}

\usepackage[margin=1.15in]{geometry}
\usepackage{amsmath,amssymb,amsthm}
\usepackage{booktabs}
\usepackage{natbib}
\usepackage[hidelinks]{hyperref}
\usepackage{setspace}

\bibpunct{(}{)}{;}{a}{}{,}

% Allow extra stretch in math-heavy paragraphs rather than overflowing the margin.
\emergencystretch=2em

% Let floats occupy more of a page before being pushed to a page of their own.
\renewcommand{\topfraction}{0.9}
\renewcommand{\bottomfraction}{0.8}
\renewcommand{\textfraction}{0.07}
\renewcommand{\floatpagefraction}{0.75}

\newtheorem{theorem}{Theorem}
\newtheorem{proposition}{Proposition}
\newtheorem{lemma}{Lemma}
\newtheorem{corollary}{Corollary}
\theoremstyle{definition}
\newtheorem{assumption}{Assumption}
\newtheorem{definition}{Definition}

\numberwithin{equation}{section}

\title{Exact Randomization Inference for Cross-Outcome\\ Falsification Tests with Few Clusters}

\author{Ethan Xie\thanks{Charter School of Wilmington, Wilmington, DE.
Email: \texttt{ethanxie@udel.edu}. Replication code and all output files
underlying the tables are available at
\url{https://github.com/opxie1/Exact-Randomization-Inference-for-Cross-Outcome-Falsification-Tests-with-Few-Clusters};
every number reported below names the file it comes from.}}

\date{August 21, 2026}

\begin{document}
\maketitle

\begin{abstract}
\noindent
Applied researchers test identifying assumptions by asking whether an estimated
treatment effect varies across outcome categories that differ in how susceptible
they are to a suspected violation. We show that the cluster-robust Wald statistic
for the joint null of no cross-category gradient is exactly a scalar multiple of
a Hotelling $T^2$, so its chi-square calibration is wrong by an amount governed
by the ratio of restrictions to clusters. At a design with $51$ clusters and $13$
post-treatment restrictions the nominal $5\%$ test rejects a true null $38.9\%$
of the time. We give three correctly sized alternatives: a label-permutation test
that is exact in finite samples under category exchangeability, a cluster
sign-change test, and a restricted wild cluster bootstrap. Permutation has a
resolution floor of $2/C!$ that makes it unusable below five categories. In a
crime-reporting application the chi-square test rejects at $p=0.018$ and no
correctly sized procedure rejects.
\end{abstract}

\smallskip
\noindent\textbf{Keywords:} randomization inference; cluster-robust inference;
falsification tests; permutation tests; Hotelling's $T^2$; difference-in-differences.

\smallskip
\noindent\textbf{JEL codes:} C12, C15, C23.

\newpage
\onehalfspacing

\section{Introduction}

Falsification tests built from several outcomes are common in applied work, and
their inference is usually asymptotic. The design is familiar: a researcher
suspects that an estimated effect is contaminated by some mechanism, observes
several outcome categories that differ in exposure to that mechanism, and tests
whether the estimated effects line up with exposure. \citet{weisburd2024crs}
run a version of this across crime categories graded by reporting propensity.
The test is then a joint restriction on interaction coefficients, and the
$p$-value comes from a chi-square reference distribution with cluster-robust
standard errors.

That calibration fails in the sample sizes where these tests are run. We derive
the exact null distribution of the cluster-robust Wald statistic for this
problem and show it is $\omega_2$ times a Hotelling $T^2$ with dimension $m$ and
$G-1$ degrees of freedom, where $m$ is the number of restrictions, $G$ is the
number of clusters, and $\omega_2$ is a fitting correction within $0.002$ of one
in our designs \citep{hotelling1931}. The chi-square reference ignores the
estimation of the $m \times m$ meat matrix from $G$ cluster scores. The
resulting size distortion is not a small-sample curiosity. At $G = 51$ and
$m = 13$ we measure a null rejection rate of $0.389$ against a nominal $0.05$,
and at $G = 20$ the rate is $0.905$.

The distortion is governed by $m$ and $G$ jointly, and the chi-square
calibration is valid only when $m$ grows slower than $G^{2/3}$. Along sequences
with $m^{3/2}/G$ converging to a positive constant the size converges to a limit
strictly above the nominal level, and when $m/G$ converges to a positive
fraction the size converges to one. This gives a design rule computable before
any outcome is observed: at $36.4$ effective clusters, size stays below $0.075$
only for $m \le 2$.

We give three procedures that are correctly sized. The first permutes the
exposure labels across outcome categories. Under an exchangeability condition on
the joint law of the category data arrays, this test has size at most $\alpha$
in finite samples for every $N$, $T$, $G$ and $C$, with no distributional,
clustering or cluster-count assumption. The condition is not free: with $C$
categories the smallest attainable $p$-value for a two-sided quadratic statistic
is $2/C!$, so the test cannot reject at the $5\%$ level with fewer than five
categories. The second procedure randomizes cluster-level sign changes and the
third is the restricted wild cluster bootstrap of \citet{cameron2008bootstrap};
both remain usable when $C$ is small.

Two results concern power rather than size. The Wald statistic is not optimal
against the alternative these tests are designed to detect. When the shape of
the response across horizons is known up to scale, a directed statistic is
uniformly most powerful invariant, and recalibrating both tests to deliver size
$\alpha$ puts the ratio of minimum detectable responses at $0.393$ at $36.4$
effective clusters. Label randomization also carries a price, which we quantify:
relative to a normal benchmark the permutation test needs $1.35$ times the
signal at $C = 5$ and $1.15$ at $C = 8$.

Section \ref{sec:setup} states the design and the null. Section \ref{sec:anatomy}
derives the exact null law of the Wald statistic and the condition under which
its chi-square calibration is valid. Section \ref{sec:perm} gives the
label-permutation test, its exactness condition, and its resolution floor.
Section \ref{sec:cluster} gives the two cluster-level procedures. Section
\ref{sec:power} derives local power, a closed-form minimum detectable effect,
and the optimal statistic. Section \ref{sec:sim} reports size and power over
$G \in \{20,30,51,100\}$ and $C \in \{3,5,8\}$. Section \ref{sec:app} applies
the procedures to two panels. Appendix \ref{app:proofs} contains the proofs.

\section{Related work}

Our exact null law is a statement about the cluster-robust Wald statistic with
$G$ clusters, which places it in the few-cluster inference literature.
\citet{cameron2008bootstrap} propose the wild cluster bootstrap with the null
imposed and document its size improvement over the cluster-robust $t$;
\citet{mackinnon2017wild} characterize when it fails. Neither paper treats the
joint $m$-restriction case in which $m$ is comparable to $G$, which is where
cross-outcome falsification tests sit. \citet{carter2017cluster} show that
cluster heterogeneity reduces the number of clusters that the asymptotics
effectively deliver, and define the effective cluster count we use in Section
\ref{sec:anatomy}. \citet{canay2017randomization} give asymptotic validity for
randomization tests when the invariance holds only in the limit, and we use
their framework for the case where our finite-sample exchangeability condition
fails. The finite-sample argument itself is the classical group-invariance
construction, for which we follow \citet{lehmann2005testing}.

The design that motivates the problem is the use of graded outcome categories to
detect a treatment-induced change in how outcomes are recorded.
\citet{weisburd2024crs} run this contrast across crime categories.
\citet{hausman1998misclassification} formalize outcome misclassification that
depends on treatment status in a cross-sectional discrete-response setting, and
\citet{denteh2026misclassdid} and \citet{gutknecht2026cic} study misclassified
outcomes in difference-in-differences and changes-in-changes designs
respectively. \citet{fhs2019pretrends} use an auxiliary series loading on a
common latent factor to purge event-study coefficients, which is the closest
existing device to the loading normalization our statistic requires. We take the
estimand and the identification argument from that literature as given. This
paper is about the reference distribution.

Falsification tests are also run as pre-trend tests, where
\citet{roth2022pretest} shows that conditioning on passing one distorts the
sampling distribution of the effect estimate, and \citet{rambachan2023credible}
replace the test with bounds on the post-treatment violation. Our results concern
the calibration of a cross-outcome test rather than the consequences of
conditioning on its outcome, and the two problems compound: a test with size
$0.39$ conditions on an event that is mostly noise.

\section{Setup and the null}
\label{sec:setup}

Units $i = 1,\dots,N$ are grouped into clusters $g = 1,\dots,G$; periods are
$t = 1,\dots,T$; outcome categories are $c = 1,\dots,C$ with known exposure
indices $d_1,\dots,d_C$. Treatment is absorbing with cohort $G_i \in
\{1,\dots,T\} \cup \{\infty\}$ and $D_{it} = \mathbf{1}\{t \ge G_i\}$.

\begin{assumption}[category-balanced design]
\label{as:B}
There is a common cell set $\mathcal{C} \subseteq \{1,\dots,N\} \times
\{1,\dots,T\}$ with $|\mathcal{C}| = n_0$ on which every category is observed,
and the treatment path $D_{it}$, the cohort $G_i$ and the cluster map $i \mapsto
g(i)$ do not depend on $c$.
\end{assumption}

Assumption \ref{as:B} rules out designs in which categories are measured on
different subsamples or at different frequencies. It is verifiable from the data
layout. Section \ref{sec:app} reports one application where it fails, and the
exact procedures below do not apply there.

Let $L$ collect the included event-time horizons with $|L| = q$, omitting
$\ell = -1$ and binning endpoints, let $X \in \mathbb{R}^{n_0 \times q}$ hold the
corresponding indicators, and let $P \subseteq \{1,\dots,q\}$ with $|P| = m$ be
the post-treatment block. With $y_{itc}$ the outcome and $\delta_c = d_c -
\bar{d}$ the centered exposure index, the pooled gradient regression is
\begin{equation}
y_{itc} = a_{ic} + p_{tc} + \sum_{\ell \in L} \beta_\ell X_{(it),\ell}
  + \sum_{\ell \in L} \gamma_\ell \delta_c X_{(it),\ell} + u_{itc},
\label{eq:pooled}
\end{equation}
estimated by OLS with the CR1 cluster-robust covariance clustered on $g$. The
hypothesis under test throughout is the joint null
\begin{equation}
H_0: \gamma_P = 0,
\label{eq:null}
\end{equation}
that the post-treatment coefficients do not vary with exposure. This is a weak
null on a parameter vector, not a sharp null of no effect for every unit, and
the distinction matters for the randomization arguments in Sections
\ref{sec:perm} and \ref{sec:cluster}: we state separately for each procedure
which invariance delivers validity and under what condition.

The cross-category structure enters \eqref{eq:pooled} as a single interaction
contrast. The procedure tests one joint hypothesis, not $C$ separate hypotheses,
so no multiplicity adjustment arises and the error rate controlled is the
per-comparison rate for \eqref{eq:null}.

Write $M$ for the two-way within projection on $\mathcal{C}$, $\tilde{X} = MX$,
$S = \tilde{X}'\tilde{X}$, and $b_c = S^{-1}\tilde{X}'M y_c$ for the
within-category event-study coefficient vector.

\begin{lemma}[separability]
\label{lem:sep}
Under Assumption \ref{as:B} and $\operatorname{rank}(\tilde{X}) = q$, the OLS
estimators in \eqref{eq:pooled} satisfy $\hat{\beta} = C^{-1}\sum_c b_c$ and
$\hat{\gamma} = \sum_c \delta_c b_c / \sum_c \delta_c^2$.
\end{lemma}

The gradient estimator is a fixed linear combination of the per-category event
studies, which is what makes the category-level randomization of Section
\ref{sec:perm} available.

The $b_c$ are within-category event studies estimated by OLS, which under
staggered adoption with heterogeneous effects need not recover a convex
combination of cohort-specific effects
\citep{dechaisemartin2020twoway,sunabraham2021}. Replacing each $b_c$ by a
heterogeneity-robust estimator \citep{callaway2021did,borusyak2024revisiting},
or by another of the alternatives surveyed in \citet{roth2023trending}, leaves
the randomization arguments of Sections \ref{sec:perm} and \ref{sec:cluster}
intact, because those require only that the category-level statistic be
order-invariant in the sense of Theorem \ref{thm:exact}. The exact null law of
Section \ref{sec:anatomy} is specific to the CR1 Wald statistic from
\eqref{eq:pooled} and would have to be re-derived for another estimator.

Failing to reject \eqref{eq:null} does not establish that the identifying
assumption holds. Section \ref{sec:power} gives the minimum detectable response
for each design, and the applications in Section \ref{sec:app} report it
alongside every non-rejection.

\section{The exact null law of the cluster-robust Wald statistic}
\label{sec:anatomy}

The size distortion has a closed form. We isolate it under an idealization that
is used only for this diagnosis and is not maintained elsewhere in the paper.

\begin{assumption}[ideal cluster experiment]
\label{as:I}
(i) Under $H_0$ in \eqref{eq:null}, the cluster score contributions
$\eta_1,\dots,\eta_G$ are i.i.d.\ $N(0,\Sigma)$ with $\Sigma_{PP} \succ 0$.
(ii) Every cluster carries the same design weight, $S_g = S/G$ for all $g$,
where $S_g = \tilde{X}_g'\tilde{X}_g$.
\end{assumption}

Part (i) holds when the errors in \eqref{eq:pooled} are Gaussian, independent
across clusters, and identically distributed across clusters given the design.
Part (ii) fails in a staggered design with a large never-treated share;
Definition \ref{def:geff} handles that case.

\begin{proposition}[exact null law]
\label{prop:anatomy}
Let Assumptions \ref{as:B} and \ref{as:I} hold, let
$\operatorname{rank}(\tilde{X}) = q$, let $H_0$ in \eqref{eq:null} hold, and let
$G - m \ge 1$. Put $\omega_2 = (n-2q)/(n-1)$ with $n = C n_0$, let $\hat{S}_P$
be the sample covariance of the $\eta_{gP}$ and $T^2 = G \bar{\eta}_P'
\hat{S}_P^{-1} \bar{\eta}_P$. Then the CR1 Wald statistic $W$ for
\eqref{eq:null} satisfies $W = \omega_2 T^2$ identically, and
\begin{equation}
\frac{(G-m)\,W}{\omega_2 (G-1) m} \sim F(m, G-m)
\label{eq:Fexact}
\end{equation}
exactly for every $G$ and $m$ with $G - m \ge 1$.
\end{proposition}

The chi-square calibration treats the estimated meat matrix as known. Under
Assumption \ref{as:I} the cost of that substitution is exactly the $F$
denominator in \eqref{eq:Fexact}, which has $G - m$ degrees of freedom. The
statement is a degrees-of-freedom result: nothing about the errors beyond
Assumption \ref{as:I}(i) enters, and $\omega_2$ equals $0.9981$ at the design of
Section \ref{sec:sim}. Proposition \ref{prop:anatomy} also supplies a valid
reference distribution, since \eqref{eq:Fexact} can be inverted for a critical
value.

\begin{corollary}[when the chi-square calibration is valid]
\label{cor:asym}
Let $\mathrm{size}(G,m)$ denote the null rejection probability of the nominal
level-$\alpha$ chi-square test under the conditions of Proposition
\ref{prop:anatomy}, and assume $\sqrt{m}\,|\log \omega_2| \to 0$.

\noindent (i) For fixed $m$, $\mathrm{size}(G,m) = \alpha + c_m f_m(c_m)
(m + c_m)/\{2(G-1)\} + O(G^{-2})$ where $c_m = \chi^2_{m,1-\alpha}$ and $f_m$ is
the $\chi^2_m$ density, so $\mathrm{size}(G,m) \to \alpha$ as $G \to \infty$.

\noindent (ii) Let $m \to \infty$ with $m/G \to 0$ and $\chi_{m,G} =
\sqrt{m/2}\,\log\{(G-1)/(G-m)\} \to \chi \in [0,\infty]$. Then
$\mathrm{size}(G,m) \to \Phi(\chi - z_{1-\alpha})$, and $\chi_{m,G} =
m^{3/2}/(\sqrt{2}\,G) + o(1)$ whenever $m^{5/2}/G^2 \to 0$. If instead
$m/G \to r \in (0,1)$, then $\mathrm{size}(G,m) \to 1$.
\end{corollary}

Part (ii) is the operative statement. The chi-square calibration is
asymptotically valid if and only if $m^{3/2}/G \to 0$, that is $m = o(G^{2/3})$.
The number of restrictions and the number of clusters cannot be considered
separately.

\begin{definition}[effective clusters]
\label{def:geff}
For a contrast $w \in \mathbb{R}^m$ put $a = S^{-1}[\cdot,P]w$ and $\varphi_g =
a'S_g a \ge 0$. Then $G_{\mathrm{eff}}(w) = (\sum_g \varphi_g)^2 / \sum_g
\varphi_g^2$.
\end{definition}

Definition \ref{def:geff} is the effective cluster count of
\citet{carter2017cluster}, specialized to the contrast $w$. It satisfies
$1 \le G_{\mathrm{eff}}(w) \le G$, with equality on the right when the
$\varphi_g$ are equal, which is what Assumption \ref{as:I}(ii) imposes.
Substituting $G_{\mathrm{eff}}$ for $G$ in \eqref{eq:Fexact} is an
approximation; Proposition \ref{prop:anatomy} does not license it. Section
\ref{sec:sim} reports its accuracy. Both $G_{\mathrm{eff}}$ and $m$ are
computable from the design alone, before any outcome is observed, which turns
Corollary \ref{cor:asym} into an ex-ante restriction on how many horizons a
researcher may test.

\section{The label-permutation test}
\label{sec:perm}

Write $\mathcal{Y}_c \in \mathbb{R}^{n_0}$ for the category-$c$ data array,
$\mathcal{D} = (\mathcal{Y}_1,\dots,\mathcal{Y}_C)$, and $\mathcal{Z}$ for the
design, which is common to all categories under Assumption \ref{as:B}. Let
$\mathcal{T}(\mathcal{D},\delta)$ be a statistic with large values evidence
against $H_0$ in \eqref{eq:null}.

\begin{definition}[label-permutation test]
\label{def:perm}
For $\pi \in S_C$ let $\delta_\pi = (\delta_{\pi(1)},\dots,\delta_{\pi(C)})$ and
$T(\pi) = \mathcal{T}(\mathcal{D},\delta_\pi)$. With $M = C!$ the test rejects
at level $\alpha$ when $\hat{p} = M^{-1}\#\{\pi \in S_C: T(\pi) \ge
T(\mathrm{id})\} \le \alpha$.
\end{definition}

The randomization is over the $C!$ assignments of the exposure labels to the
category data arrays, holding the arrays and the design fixed. Nothing is
permuted at the unit, cluster or period level.

\begin{assumption}[category exchangeability]
\label{as:CE}
Conditional on $\mathcal{Z}$, $(\mathcal{Y}_{\pi(1)},\dots,
\mathcal{Y}_{\pi(C)}) \overset{d}{=} (\mathcal{Y}_1,\dots,\mathcal{Y}_C)$ for
every $\pi \in S_C$, and $d_1,\dots,d_C$ are fixed constants.
\end{assumption}

Assumption \ref{as:CE} rules out categories with different error variances,
different loadings on a common factor unless the outcome is normalized, and any
dependence of the labels on the data. It is not implied by $H_0$ in
\eqref{eq:null} and it is not directly testable, though its leading failure mode
has an observable signature: Section \ref{sec:sim} prices the size distortion as
a function of the ratio of the largest to smallest category error variance.

\begin{theorem}[exact finite-sample validity]
\label{thm:exact}
Suppose Assumption \ref{as:CE} holds and $\mathcal{T}$ is order-invariant, in
the sense that it is a function of the multiset
$\{(\delta_c,\mathcal{Y}_c)\}_{c=1}^C$ and not of the order in which the pairs
are listed. Then for every $N$, $T$, $G$ and $C$, the test of Definition
\ref{def:perm} for the null \eqref{eq:null} under the randomization over $S_C$
satisfies $\Pr(\text{reject}) \le \alpha$, with $\Pr(\text{reject}) =
\lfloor \alpha M \rfloor / M$ when the $M$ values are almost surely distinct.
The randomized version has size exactly $\alpha$.
\end{theorem}

Exactness here is finite-sample exactness for the null \eqref{eq:null} under the
randomization distribution of Definition \ref{def:perm}, and it requires
Assumption \ref{as:CE}. No condition on $G$, on the error distribution, or on
the within-cluster dependence structure is used. Order-invariance holds for the
OLS estimates and CR1 covariance in \eqref{eq:pooled}, and it fails for any
loading normalization that selects a reference category by its position in an
ordering of $d$.

\begin{proposition}[attainable $p$-values]
\label{prop:resolution}
(i) $\hat{p} \ge 1/C!$ always. (ii) If $\mathcal{T}$ is even in the labels, so
that $\mathcal{T}(\mathcal{D},-\delta) = \mathcal{T}(\mathcal{D},\delta)$, and
the multiset $\{\delta_c\}$ is symmetric about zero, then $\hat{p} \ge 2/C!$.
(iii) A level-$\alpha$ label-permutation test therefore requires $C! \ge
2/\alpha$ for the two-sided Wald statistic with equally spaced $d$, that is
$C \ge 5$ at $\alpha = 0.05$, and $C! \ge 1/\alpha$ for a one-sided statistic,
that is $C \ge 4$.
\end{proposition}

The Wald statistic is even in the labels because $\hat{\gamma}(-\delta) =
-\hat{\gamma}(\delta)$ and its covariance estimator is quadratic in the cluster
scores. Proposition \ref{prop:resolution} is a property of the randomization
group, not of any implementation. With $C = 3$ the two-sided Wald has three
attainable $p$-values, the smallest being $1/3$, and the test can never reject
at conventional levels. Section \ref{sec:sim} reports the resulting rejection
rate of exactly $0.000$.

\section{Cluster-level randomization and the bootstrap}
\label{sec:cluster}

Two procedures remain available when $C \le 4$. Both randomize over clusters
rather than category labels, so Proposition \ref{prop:resolution} does not bind.

The first randomizes sign changes over $\mathcal{E} = \{-1,+1\}^G$, applied to
the cluster contributions $\zeta_1,\dots,\zeta_G$ to the gradient score with the
null \eqref{eq:null} imposed.

\begin{theorem}[cluster sign-change test]
\label{thm:signflip}
If, under $H_0$ in \eqref{eq:null}, the cluster contributions are jointly
sign-symmetric, meaning $(\epsilon_1 \zeta_1,\dots,\epsilon_G \zeta_G)
\overset{d}{=} (\zeta_1,\dots,\zeta_G)$ for every $\epsilon \in \mathcal{E}$,
then the sign-change test is exactly level $\alpha$ in finite samples, and so is
its Monte Carlo version using $B-1$ random sign vectors together with the
identity. Sign symmetry holds for the infeasible statistic built from the true
errors under \eqref{eq:null} whenever the cluster error contributions are
independent across $g$ and symmetric about zero. The feasible statistic differs
from the infeasible one by $O_p(G^{-1/2})$ and is asymptotically level $\alpha$
as $G \to \infty$.
\end{theorem}

Exactness for the feasible version does not follow from Theorem
\ref{thm:signflip}, and we do not claim it. The finite-sample statement applies
to the infeasible statistic; the feasible statistic inherits validity only
asymptotically in $G$, by the approximate-symmetry argument of
\citet{canay2017randomization}. Section \ref{sec:sim} measures the finite-sample
gap.

The second procedure is the restricted wild cluster bootstrap of
\citet{cameron2008bootstrap}, drawing Rademacher weights at the cluster level
with \eqref{eq:null} imposed on the residuals. \citet{mackinnon2017wild}
document the conditions under which it fails, and those conditions concern
cluster heterogeneity rather than the ratio $m/G$ that drives Corollary
\ref{cor:asym}.

\section{Power, the minimum detectable response, and the optimal statistic}
\label{sec:power}

Let the alternative be $\gamma_\ell = \kappa\, k(\ell)$ for a response path
$k(\cdot) \ge 0$ and a scalar magnitude $\kappa > 0$, so that $H_1: \kappa > 0$.
Under local alternatives the noncentrality of the Wald statistic depends on the
design through two scalars: $\Sigma_\delta^2 = \sum_c (d_c - \bar{d})^2$, a
property of the exposure index, and a timing quantity $I_k$ determined by the
treatment-timing design and the horizon window.

\begin{corollary}[closed-form minimum detectable response]
\label{cor:mde}
Let $\lambda^*(\alpha,\beta,m)$ solve $\Pr\{\chi^2_m(\lambda^*) >
\chi^2_{m,1-\alpha}\} = 1-\beta$. With a linear wedge and homoskedastic
categories,
\begin{equation}
\kappa^{\mathrm{Wald}}_{\mathrm{MDE}} = \sigma
\sqrt{\frac{\lambda^*(\alpha,\beta,m)}{\Sigma_\delta^2\, I_k}},
\qquad
\kappa^{\mathrm{dir}}_{\mathrm{MDE}} = \sigma\,
\frac{z_{1-\alpha}+z_\beta}{\sqrt{\Sigma_\delta^2\, I_k}} .
\label{eq:mde}
\end{equation}
\end{corollary}

Both $\Sigma_\delta^2$ and $I_k$ are computable before any outcome is observed,
and $\sigma^2$ is the only quantity requiring pilot data, so \eqref{eq:mde} is
an ex-ante design rule. The expression for $\kappa^{\mathrm{Wald}}_{\mathrm{MDE}}$
is stated against the chi-square critical value, which Proposition
\ref{prop:anatomy} shows is the wrong one; it therefore understates the minimum
detectable response of any correctly calibrated Wald test. Replacing
$\chi^2_{m,1-\alpha}$ by the exact critical value from \eqref{eq:Fexact} raises
the required noncentrality from $17.85$ to $41.46$ at $G_{\mathrm{eff}} = 36.42$
and $m = 13$, inflating the minimum detectable response by a factor of $1.524$.

The Wald statistic is not the right statistic against this alternative.

\begin{theorem}[optimal directed test]
\label{thm:umpi}
Suppose $V = \operatorname{Var}(\hat{\gamma}_P)$ and the response shape $k_P$
are known. Among tests of $H_0: \kappa = 0$ against $H_1: \kappa > 0$ that are
invariant under shifts of the common event-study component and depend on the
data only through the maximal invariant, the test rejecting when $t = k_P'V^{-1}
\hat{\gamma}_P / \sqrt{k_P'V^{-1}k_P} > z_{1-\alpha}$ is uniformly most
powerful, with power $\Phi(\kappa \sqrt{k_P'V^{-1}k_P} - z_{1-\alpha})$.
\end{theorem}

When $k_P$ is restricted only to the non-negative orthant no uniformly most
powerful invariant test exists, and the likelihood-ratio statistic has the
chi-bar-square null distribution of \citet{kudo1963} and \citet{shapiro1988},
a mixture $\sum_{j=0}^m w_j(V)\chi^2_j$ in which $w_j(V)$ is the probability
that the projection of $N(0,V)$ onto the non-negative orthant in the $V^{-1}$
metric has exactly $j$ positive coordinates. Its power lies between the directed
test and the two-sided Wald.

Calibrating each test so that it delivers size $\alpha$, the ratio of minimum
detectable responses between the directed and Wald tests is $0.393$ at
$G_{\mathrm{eff}} = 36.42$, $0.448$ at $G = 51$ and $0.516$ at $G = 100$.
Directing the test cuts the detectable response by roughly $61\%$ at the
smallest of these designs.

\section{Simulations}
\label{sec:sim}

\paragraph{Design.} We simulate a two-layer panel. A latent factor $F_{it} =
f_i + g_t + \tau D_{it} + v_{it}$ drives $C$ category outcomes $Y^*_{itc} =
a_{ic} + \lambda_{tc} + \theta_c F_{it} + \varepsilon_{itc}$, and the recorded
outcome adds a time-invariant unit-by-category distortion and a
treatment-induced term, $Y_{itc} = Y^*_{itc} + \mu_{ic} + \kappa k(\ell_{it})
d_c D_{it} + u_{itc}$. All disturbances are Gaussian and independent across
clusters. The null \eqref{eq:null} corresponds to $\kappa = 0$. We vary
$G \in \{20,30,51,100\}$ with one unit per cluster, $C \in \{3,5,8\}$, and set
$T = 120$ with three adoption cohorts, $q = 18$ included horizons and $m = 13$
post-treatment restrictions. Exposure indices are equally spaced. Every size
cell uses $2{,}000$ replications at nominal level $0.05$; permutation tests
enumerate all $C!$ label assignments and the bootstrap and sign-change tests use
$1{,}999$ draws. Monte Carlo standard errors for a rejection rate near $0.05$
are $0.005$.

\paragraph{Size.} Table \ref{tab:size} reports null rejection rates. The
chi-square calibration rejects at $0.920$, $0.693$, $0.393$ and $0.172$ at
$G = 20$, $30$, $51$ and $100$ with $C = 3$, and the pattern is nearly identical
across $C$, consistent with Proposition \ref{prop:anatomy}, in which $C$ enters
only through $\omega_2$. The $F$ calibration reduces but does not remove the
distortion. The permutation test is at nominal size for $C \in \{5,8\}$,
ranging from $0.0425$ to $0.059$ across all eight such cells. At $C = 3$ it
rejects exactly $0.000$ of the time, which is Proposition
\ref{prop:resolution}(iii) rather than a numerical artifact. The wild cluster
bootstrap ranges from $0.0505$ to $0.067$ and the sign-change test from $0.0215$
to $0.066$, the latter conservative at $G = 20$.

\paragraph{Accuracy of the exact law.} Evaluating \eqref{eq:Fexact} at
$G_{\mathrm{eff}}$ from Definition \ref{def:geff} predicts a null rejection rate
of $0.378$ against the observed $0.389$ at $G = 51$, $C = 5$, $m = 13$. Across
the twelve $(G,C)$ cells the mean absolute prediction error is $2.37$ percentage
points using $G_{\mathrm{eff}}$ and $13.13$ percentage points using $G$. In this
design $G_{\mathrm{eff}} = 36.42$, or $71.4\%$ of $G$, and it is a design object
rather than a random one: its standard deviation across replications is
$10^{-13}$. The exact Hotelling calibration evaluated at $G_{\mathrm{eff}}$
rejects at $0.0445$ to $0.0470$ at $G = 51$ and is conservative at $G = 20$,
where it rejects at $0.000$.

\paragraph{The design rule.} Applying Corollary \ref{cor:asym} through
\eqref{eq:Fexact}, the largest number of restrictions holding chi-square size
below $0.075$ is $m = 2$ at $G_{\mathrm{eff}} = 36.4$, $m = 4$ at $72.7$, $m =
6$ at $100$ and $m = 8$ at $150$. Testing $m = 13$ restrictions at
$G_{\mathrm{eff}} = 36.4$ gives a predicted size of $0.379$.

\paragraph{When exchangeability fails.} We vary the ratio of the largest to
smallest category error variance, which violates Assumption \ref{as:CE}. At
$C = 5$ the permutation test using the Wald statistic rejects at $0.050$,
$0.0635$, $0.0772$ and $0.0901$ for variance ratios $1$, $2$, $4$ and $9$. The
distortion is monotone in the ratio and is an order of magnitude smaller than
the chi-square distortion it replaces.

\paragraph{The price of label randomization.} Relative to a normal benchmark,
the permutation test requires a noncentrality inflated by $1.892$ at $C = 4$,
$1.352$ at $C = 5$, $1.235$ at $C = 6$, $1.149$ at $C = 8$ and $1.098$ at
$C = 10$. Exactness is bought with power, and the price falls as $C$ grows.

\paragraph{Power.} Table \ref{tab:power} reports size-adjusted power at
$G = 51$ against a linear recording response of magnitude $\kappa$, over $600$
replications per cell. Two patterns are visible and they point in opposite
directions. Label randomization with the two-sided Wald statistic is the least
powerful procedure at every $\kappa$ and every $C$ we ran: at $C = 5$ and
$\kappa = 0.45$ it attains $0.383$ against $0.705$ for the chi-square statistic
at its size-adjusted critical value. Directing the statistic more than reverses
that loss. The directed cluster sign-change test attains $0.983$ in the same
cell, and at $C = 8$ it reaches $0.842$ at $\kappa = 0.20$ where the two-sided
Wald reaches $0.277$. A researcher who buys exactness with label permutation and
then spends it on a two-sided statistic is worse off than before; the gain comes
from the direction restriction, not from the randomization.

The size-adjusted power of the two-sided permutation test is non-monotone in
$\kappa$ at $C = 5$, falling from $0.452$ at $\kappa = 0.60$ to $0.383$ at
$\kappa = 0.80$. The Monte Carlo standard error at these rates is $0.020$, so
the decline is about three standard errors and we do not attribute it to
simulation noise. We have not established its cause.

\begin{table}[htbp]
\centering
\caption{Null rejection rates, nominal level $0.05$}
\label{tab:size}
\begin{tabular}{llrrrr}
\toprule
 & & \multicolumn{4}{c}{Clusters $G$} \\
\cmidrule(lr){3-6}
Calibration & $C$ & 20 & 30 & 51 & 100 \\
\midrule
Chi-square             & 3 & 0.920 & 0.693 & 0.393 & 0.172 \\
                       & 5 & 0.905 & 0.686 & 0.389 & 0.187 \\
                       & 8 & 0.917 & 0.690 & 0.356 & 0.172 \\
\addlinespace
$F$                    & 3 & 0.838 & 0.586 & 0.311 & 0.143 \\
                       & 5 & 0.821 & 0.558 & 0.309 & 0.153 \\
                       & 8 & 0.834 & 0.571 & 0.283 & 0.138 \\
\addlinespace
Label permutation      & 3 & 0.000 & 0.000 & 0.000 & 0.000 \\
                       & 5 & 0.053 & 0.050 & 0.046 & 0.059 \\
                       & 8 & 0.051 & 0.051 & 0.043 & 0.058 \\
\addlinespace
Wild cluster bootstrap & 3 & 0.053 & 0.055 & 0.053 & 0.061 \\
                       & 5 & 0.064 & 0.062 & 0.051 & 0.067 \\
                       & 8 & 0.061 & 0.061 & 0.051 & 0.056 \\
\addlinespace
Cluster sign-change    & 3 & 0.027 & 0.036 & 0.046 & 0.058 \\
                       & 5 & 0.022 & 0.041 & 0.050 & 0.066 \\
                       & 8 & 0.033 & 0.041 & 0.048 & 0.053 \\
\bottomrule
\end{tabular}

\smallskip
\begin{minipage}{0.95\textwidth}
\footnotesize Note: rejection rates for the joint null $H_0: \gamma_P = 0$ of
\eqref{eq:null} in the pooled gradient regression \eqref{eq:pooled}, estimated
by OLS with the CR1 covariance clustered on $g$. Design: $m = 13$
post-treatment restrictions, $q = 18$ included horizons, $T = 120$, three
adoption cohorts, one unit per cluster, equally spaced exposure indices. Label
permutation randomizes the $C!$ assignments of the exposure index to category
data arrays and enumerates all of them; the wild cluster bootstrap and the
cluster sign-change test use $1{,}999$ draws. $2{,}000$ replications per cell;
Monte Carlo standard error at a rate of $0.05$ is $0.005$. Source:
\texttt{results/theory/inference/size\_table.csv}.
\end{minipage}
\end{table}

\begin{table}[htbp]
\centering
\caption{Size-adjusted power against the recording response $\kappa$, $G = 51$}
\label{tab:power}
\begin{tabular}{llrrrrrr}
\toprule
 & & \multicolumn{6}{c}{Recording response $\kappa$} \\
\cmidrule(lr){3-8}
Statistic and calibration & $C$ & 0.05 & 0.10 & 0.15 & 0.20 & 0.30 & 0.45 \\
\midrule
Wald, chi-square              & 5 & 0.082 & 0.107 & 0.145 & 0.212 & 0.362 & 0.705 \\
                              & 8 & 0.058 & 0.113 & 0.150 & 0.277 & 0.537 & 0.850 \\
\addlinespace
Wald, wild cluster bootstrap  & 5 & 0.080 & 0.110 & 0.147 & 0.220 & 0.362 & 0.715 \\
                              & 8 & 0.053 & 0.112 & 0.157 & 0.273 & 0.532 & 0.852 \\
\addlinespace
Wald, label permutation       & 5 & 0.083 & 0.103 & 0.138 & 0.170 & 0.212 & 0.383 \\
                              & 8 & 0.055 & 0.097 & 0.142 & 0.187 & 0.395 & 0.655 \\
\addlinespace
Directed, label permutation   & 5 & 0.147 & 0.250 & 0.403 & 0.518 & 0.657 & 0.797 \\
                              & 8 & 0.175 & 0.370 & 0.535 & 0.732 & 0.890 & 0.973 \\
\addlinespace
Directed, cluster sign-change & 5 & 0.143 & 0.320 & 0.527 & 0.688 & 0.895 & 0.983 \\
                              & 8 & 0.155 & 0.403 & 0.613 & 0.842 & 0.968 & 0.997 \\
\bottomrule
\end{tabular}

\smallskip
\begin{minipage}{0.95\textwidth}
\footnotesize Note: size-adjusted rejection rates against the alternative
$\gamma_\ell = \kappa k(\ell)$ with a linear response path, for the joint null
$H_0: \gamma_P = 0$ of \eqref{eq:null}. Each entry uses the critical value that
delivers rejection rate $0.05$ for that procedure under $\kappa = 0$ in the same
design, so entries are comparable across rows. Design: $G = 51$ clusters,
$m = 13$ post-treatment restrictions, $q = 18$ horizons, $T = 120$, three
adoption cohorts, equally spaced exposure indices. Label permutation enumerates
all $C!$ assignments; the cluster sign-change test uses $1{,}999$ draws. $600$
replications per cell; Monte Carlo standard error at a rate of $0.5$ is $0.020$.
Source: \texttt{results/theory/inference/power\_curves.csv}.
\end{minipage}
\end{table}

\section{Application}
\label{sec:app}

We apply the procedures to two state panels assembled for a companion project on
treatment-induced recording change. The identification claim in each case is
that under measurement invariance the post-treatment event-study coefficients do
not vary with the exposure index, so \eqref{eq:null} holds. The application is
illustrative of the inference problem rather than a contribution to either
substantive literature.

The first panel covers $51$ states from 1979 to 2024 with six offense categories
graded by the share of victimizations reported to police, clustered on state, so
$G = 51$, $C = 6$ and $m = 9$. The design has $G_{\mathrm{eff}} = 45.9$, which
by the rule of Section \ref{sec:sim} supports $m = 3$ restrictions at size
$0.075$; the specification tests $m = 9$, for which \eqref{eq:Fexact} predicts a
chi-square size of $0.169$. The chi-square test of \eqref{eq:null} rejects, with
a statistic of $19.91$ and $p = 0.018$. No correctly sized procedure rejects:
the exact Hotelling calibration at $G_{\mathrm{eff}}$ returns $p = 0.097$, the
label-permutation test over all $720$ assignments returns $p = 0.096$, the wild
cluster bootstrap returns $p = 0.075$ and the cluster sign-change test returns
$p = 0.083$. The reported finding is an artifact of the reference distribution.

A pre-period placebo in the same specification behaves differently. Its
chi-square $p$-value is $9 \times 10^{-6}$ and its label-permutation $p$-value
is $0.0167$, which is the smallest attainable value at $C = 6$ under
Proposition \ref{prop:resolution}(ii). The placebo rejection survives correct
calibration, so the pre-period contamination in that design is not an inference
artifact.

The second panel covers $51$ states from 2011 to 2024 with health conditions
graded by whether a clinical encounter is required for the respondent to report
them. Two conditions are surveyed only in odd years, so Assumption \ref{as:B}
fails on the full category set and the exact procedures of Sections
\ref{sec:perm} and \ref{sec:cluster} do not apply to the specification as
originally run. We report two balanced sub-designs instead: five annually
surveyed categories, and all seven categories on odd years only. Neither rejects
\eqref{eq:null} under any calibration. Both have $G_{\mathrm{eff}}$ near $33$
and support $m = 2$ restrictions against the $m = 7$ tested.

Neither non-rejection is evidence that measurement invariance holds. The
minimum detectable response from \eqref{eq:mde} exceeds the magnitude implied by
the substantive mechanism in both panels, and we report that comparison rather
than describing the tests as passed.

\section{Conclusion}

The cluster-robust Wald statistic for a cross-outcome falsification test is a
scaled Hotelling $T^2$, so its chi-square calibration is valid only when the
number of restrictions grows more slowly than the two-thirds power of the number
of clusters. At $51$ clusters and $13$ restrictions the nominal $5\%$ test
rejects a true null $38.9\%$ of the time. Three procedures restore size: an
exact label-permutation test valid in finite samples under category
exchangeability, a cluster sign-change test, and the restricted wild cluster
bootstrap. In a crime-reporting design the chi-square test rejects at $p = 0.018$
and none of the three does.

The main limitation is Assumption \ref{as:CE}. Exactness of the
label-permutation test requires the category data arrays to be exchangeable
given the design, which fails when categories have different error variances,
and our simulations put the size at $0.090$ when the largest category variance
is nine times the smallest. The cluster-level procedures do not require it, and
they are the default when the exchangeability condition is doubtful or when
$C \le 4$ makes permutation unusable.

One extension is immediate. Theorem \ref{thm:exact} requires a
category-balanced design, and the health application shows that requirement
binding in practice. Extending the exact reduction to panels unbalanced across
categories would recover the procedure for designs where categories are
surveyed on different schedules.

\appendix
\section{Proofs}
\label{app:proofs}

\begin{proof}[Proof of Lemma \ref{lem:sep}]
Under Assumption \ref{as:B} the regressor block in \eqref{eq:pooled} is common
across categories after the two-way within transformation, so stacking the $C$
categories gives design matrix $[\,\mathbf{1}_C \otimes \tilde{X},\ \delta
\otimes \tilde{X}\,]$ with $\delta = (\delta_1,\dots,\delta_C)'$. Because
$\mathbf{1}_C'\delta = 0$ by the centering of $\delta$, the two blocks are
orthogonal, and the normal equations separate into
$(\mathbf{1}_C'\mathbf{1}_C) S \hat{\beta} = \sum_c \tilde{X}'\tilde{y}_c$ and
$(\delta'\delta) S \hat{\gamma} = \sum_c \delta_c \tilde{X}'\tilde{y}_c$.
Dividing by $C$ and $\sum_c \delta_c^2$ respectively and using $b_c =
S^{-1}\tilde{X}'\tilde{y}_c$ gives the stated expressions.
\end{proof}

\begin{proof}[Proof of Proposition \ref{prop:anatomy}]
Write the CR1 covariance estimator as $\hat{V} = \omega_2^{-1}\,
\frac{G}{G-1}\, A \big(\sum_g a_g a_g'\big) A'$ with $A$ the relevant rows of
$S^{-1}$ and $a_g$ the cluster score. Under Assumption \ref{as:I}(ii) each
cluster carries design weight $S_g = S/G$, so $A a_g = G^{-1}\eta_g$ up to the
common factor absorbed in $A$, and the gradient estimator satisfies
$\hat{\gamma}_P = \bar{\eta}_P$ in the same normalization. Substituting,
$$
W = \hat{\gamma}_P' \hat{V}_{PP}^{-1} \hat{\gamma}_P
  = \omega_2\, \bar{\eta}_P' \Big[\tfrac{1}{G-1}\textstyle\sum_g
    (\eta_{gP}-\bar{\eta}_P)(\eta_{gP}-\bar{\eta}_P)'\Big]^{-1}
    \bar{\eta}_P \, G
  = \omega_2 T^2 ,
$$
where the centering appears because the null \eqref{eq:null} is imposed on the
residuals. This is an algebraic identity and does not use Assumption
\ref{as:I}(i).

For the distribution, Assumption \ref{as:I}(i) gives $\sqrt{G}\,\bar{\eta}_P
\sim N(0,\Sigma_{PP})$ and $(G-1)\hat{S}_P \sim W_m(\Sigma_{PP}, G-1)$
independently, by Cochran's theorem applied to the i.i.d.\ Gaussian sample
$\eta_{1P},\dots,\eta_{GP}$. Hence $T^2$ is Hotelling's $T^2$ with dimension $m$
and $G-1$ degrees of freedom \citep{hotelling1931}, and the standard
transformation $\frac{G-m}{(G-1)m}T^2 \sim F(m,G-m)$ holds for $G-m\ge1$.
Substituting $T^2 = W/\omega_2$ gives \eqref{eq:Fexact}.
\end{proof}

\begin{proof}[Proof of Corollary \ref{cor:asym}]
By \eqref{eq:Fexact} the exact size of the nominal chi-square test is
$\Pr\{F(m,G-m) > (G-m)c_m / (\omega_2(G-1)m)\}$ with $c_m =
\chi^2_{m,1-\alpha}$. Write $F(m,G-m) \overset{d}{=} (\chi^2_m/m) /
(\chi^2_{G-m}/(G-m))$ with the two chi-squares independent, so the size equals
$\Pr\{\chi^2_m > c_m \,\omega_2^{-1}(G-1)^{-1}\chi^2_{G-m}\}$.

For (i), fix $m$ and expand $\chi^2_{G-m}/(G-1) = 1 + O_p(G^{-1/2})$ with mean
$1 - (m-1)/(G-1)$. A second-order Taylor expansion of $u \mapsto
\Pr\{\chi^2_m > c_m u\}$ around $u = 1$, together with $\sqrt{m}|\log\omega_2|
\to 0$, gives the stated expansion, in which the $O(G^{-1})$ term carries
coefficient $c_m f_m(c_m)(m+c_m)/2 > 0$.

For (ii), take logarithms. The event is $\log \chi^2_m - \log \chi^2_{G-m} >
\log c_m - \log(\omega_2(G-1))$. With $m \to \infty$, $\log\chi^2_m$ is
asymptotically normal with mean $\log m - 1/m + O(m^{-2})$ and variance $2/m$,
and similarly for $\log\chi^2_{G-m}$ with $G - m \to \infty$. Standardizing and
using $\log c_m = \log m + \sqrt{2/m}\,z_{1-\alpha} + o(m^{-1/2})$ gives size
$\to \Phi(\chi - z_{1-\alpha})$ with $\chi = \lim \sqrt{m/2}\,\log\{(G-1)/(G-m)\}$.
Expanding $\log\{(G-1)/(G-m)\} = (m-1)/(G-1) + O((m/G)^2)$ gives $\chi_{m,G} =
m^{3/2}/(\sqrt{2}G) + o(1)$ when $m^{5/2}/G^2 \to 0$. If $m/G \to r \in (0,1)$
then $\log\{(G-1)/(G-m)\} \to -\log(1-r) > 0$ and $\chi_{m,G} \to \infty$, so
the size tends to one.
\end{proof}

\begin{proof}[Proof of Theorem \ref{thm:exact}]
The argument is the classical group-invariance construction; see
\citet[Ch.~15]{lehmann2005testing}. Order-invariance of $\mathcal{T}$ gives, for
every $\pi \in S_C$, $\mathcal{T}(\mathcal{D},\delta_\pi) =
\mathcal{T}(\mathcal{D}_{\pi^{-1}},\delta)$ where $\mathcal{D}_{\pi^{-1}} =
(\mathcal{Y}_{\pi^{-1}(1)},\dots,\mathcal{Y}_{\pi^{-1}(C)})$: the statistic is
computed from the multiset $\{(\delta_{\pi(c)},\mathcal{Y}_c)\}_c$, and
re-indexing by $j = \pi(c)$ turns this into
$\{(\delta_j,\mathcal{Y}_{\pi^{-1}(j)})\}_j$.

Under Assumption \ref{as:CE}, $\mathcal{D}_{\pi^{-1}} \overset{d}{=}
\mathcal{D}$ conditional on $\mathcal{Z}$ for each fixed $\pi$. Since $S_C$ is a
finite group acting on the index set, the vector $(T(\pi))_{\pi \in S_C}$ is
exchangeable conditional on $\mathcal{Z}$: for any $\sigma \in S_C$,
$(T(\sigma\pi))_{\pi} \overset{d}{=} (T(\pi))_{\pi}$. Exchangeability of the $M$
values implies that the rank of $T(\mathrm{id})$ among them is uniform on
$\{1,\dots,M\}$ when ties have probability zero, and stochastically no larger
than uniform otherwise. Rejecting when $T(\mathrm{id})$ exceeds the
$k$th order statistic with $k = \lceil (1-\alpha)M\rceil$ therefore has
conditional probability $(M-k)/M \le \alpha$, and $\lfloor \alpha M\rfloor/M$
under distinctness. Averaging over $\mathcal{Z}$ preserves the bound. The
randomized version distributes the residual mass $\alpha - \lfloor \alpha
M\rfloor/M$ over the ties and attains $\alpha$ exactly.
\end{proof}

\begin{proof}[Proof of Proposition \ref{prop:resolution}]
(i) The identity permutation contributes one count to the numerator of
$\hat{p}$, so $\hat{p} \ge 1/M$ with $M = C!$.

(ii) With equally spaced $d$ we may write $\delta_c = s\,(c - \frac{C+1}{2})$ for
a scale $s>0$, so the order-reversing permutation $r(c) = C+1-c$ satisfies
$\delta_{r(c)} = -\delta_c$. Pre-composition negates the label vector:
$(\delta_{r \circ \pi})_c = \delta_{r(\pi(c))} = -\delta_{\pi(c)}$, so
$\delta_{r\circ\pi} = -\delta_\pi$. Evenness of $\mathcal{T}$ then gives
$T(r\circ\pi) = T(\pi)$ for every $\pi$. The map $\pi \mapsto r\circ\pi$ is left
translation by an involution without fixed points, hence a fixed-point-free
involution on $S_C$, so the $M$ values fall into $M/2$ equal pairs and every
attainable $\hat{p}$ is a multiple of $2/M$.

For the Wald statistic, $\hat{\gamma}(-\delta) = -\hat{\gamma}(\delta)$ by
Lemma \ref{lem:sep}, each cluster score is odd in $\delta$, and the CR1
covariance is quadratic in the cluster scores, so $\hat{V}_\gamma(-\delta) =
\hat{V}_\gamma(\delta)$ and the quadratic form is unchanged.

(iii) Rejection at level $\alpha$ requires some attainable $\hat p \le \alpha$,
hence $2/C! \le \alpha$ in the even case and $1/C! \le \alpha$ otherwise.
\end{proof}

\begin{proof}[Proof of Theorem \ref{thm:signflip}]
For the first claim, replace $S_C$ in the proof of Theorem \ref{thm:exact} by
the finite group $\mathcal{E} = \{-1,+1\}^G$ acting by coordinatewise sign
change. That proof uses only finiteness of the group and invariance of the data
law under it, both of which the sign-symmetry hypothesis supplies.

For the second claim, let $u^0$ denote the true errors under \eqref{eq:null}.
The infeasible contribution $\zeta^0_g$ is a fixed linear functional of cluster
$g$'s errors alone, so independence across $g$ together with symmetry of each
summand about zero gives $(\epsilon_g \zeta^0_g)_g \overset{d}{=} (\zeta^0_g)_g$
for every $\epsilon \in \mathcal{E}$.

For the third claim, the feasible contribution replaces $u^0$ by the restricted
residual, which subtracts a term $S^{-1}S_g[\sum_h(\cdot)_h]$. The factor
$\|S_g\|/\|S\|$ is $O(G^{-1})$ and it multiplies a sum of $G$ cluster
contributions, which is $O_p(G^{1/2})$ times a single one; relative to $\zeta_g$
itself the discrepancy is therefore $O_p(G^{-1/2})$. Asymptotic level $\alpha$
follows from the approximate-symmetry theorem of
\citet{canay2017randomization}, whose conditions require the randomization
distribution to converge to one that is invariant under $\mathcal{E}$, which the
$O_p(G^{-1/2})$ bound delivers.
\end{proof}

\begin{proof}[Proof of Corollary \ref{cor:mde}]
Under the stated alternative and homoskedasticity, $\hat{\gamma}_P$ is
asymptotically normal with mean $\kappa k_P$ and covariance $V =
\sigma^2 (\Sigma_\delta^2 I_k)^{-1} \Omega_P$ for a design matrix $\Omega_P$
normalized so that the linear wedge gives $k_P'\Omega_P^{-1}k_P = 1$. The Wald
statistic is then noncentral chi-square with $m$ degrees of freedom and
noncentrality $\lambda = \kappa^2 \Sigma_\delta^2 I_k/\sigma^2$. Setting
$\lambda = \lambda^*(\alpha,\beta,m)$ and solving for $\kappa$ gives the first
expression in \eqref{eq:mde}. For the directed statistic, $t$ is standard normal
under \eqref{eq:null} and normal with mean $\kappa \sqrt{\Sigma_\delta^2
I_k}/\sigma$ under the alternative, so power $1-\beta$ at level $\alpha$
requires that mean to equal $z_{1-\alpha}+z_\beta$, which gives the second
expression.
\end{proof}

\begin{proof}[Proof of Theorem \ref{thm:umpi}]
Invariance under shifts of the common event-study component reduces the data to
the maximal invariant $\hat{\gamma}_P$, which is $N(\kappa k_P, V)$ with $V$ and
$k_P$ known. The family in $\kappa$ has density proportional to
$\exp\{\kappa k_P'V^{-1}\hat{\gamma}_P - \tfrac{1}{2}\kappa^2 k_P'V^{-1}k_P\}$
relative to the $\kappa = 0$ density, a one-parameter exponential family with
natural sufficient statistic $k_P'V^{-1}\hat{\gamma}_P$ and natural parameter
$\kappa$. By the Neyman-Pearson lemma the most powerful test against any fixed
$\kappa > 0$ rejects for large values of that statistic, and the rejection
region does not depend on which $\kappa > 0$ is chosen, so the test is uniformly
most powerful among invariant tests. Standardizing by
$\sqrt{k_P'V^{-1}k_P}$ gives $t \sim N(\kappa\sqrt{k_P'V^{-1}k_P},1)$, from
which the size and the stated power function follow.
\end{proof}

\begin{thebibliography}{99}
\itemsep 3pt

\bibitem[Borusyak et al.(2024)]{borusyak2024revisiting}
Borusyak, K., X.~Jaravel, and J.~Spiess (2024). Revisiting event-study designs:
Robust and efficient estimation. \emph{Review of Economic Studies} 91(6),
3253--3285.

\bibitem[Cameron et al.(2008)]{cameron2008bootstrap}
Cameron, A.~C., J.~B. Gelbach, and D.~L. Miller (2008). Bootstrap-based
improvements for inference with clustered errors. \emph{Review of Economics and
Statistics} 90(3), 414--427.

\bibitem[Canay et al.(2017)]{canay2017randomization}
Canay, I.~A., J.~P. Romano, and A.~M. Shaikh (2017). Randomization tests under
an approximate symmetry assumption. \emph{Econometrica} 85(3), 1013--1030.

\bibitem[Carter et al.(2017)]{carter2017cluster}
Carter, A.~V., K.~T. Schnepel, and D.~G. Steigerwald (2017). Asymptotic
behavior of a $t$-test robust to cluster heterogeneity. \emph{Review of
Economics and Statistics} 99(4), 698--709.

\bibitem[Callaway and Sant'Anna(2021)]{callaway2021did}
Callaway, B. and P.~H.~C. Sant'Anna (2021). Difference-in-differences with
multiple time periods. \emph{Journal of Econometrics} 225(2), 200--230.

\bibitem[de Chaisemartin and D'Haultf\oe uille(2020)]{dechaisemartin2020twoway}
de Chaisemartin, C. and X.~D'Haultf\oe uille (2020). Two-way fixed effects
estimators with heterogeneous treatment effects. \emph{American Economic
Review} 110(9), 2964--2996.

\bibitem[Denteh and K\'edagni(2022)]{denteh2026misclassdid}
Denteh, A. and D.~K\'edagni (2022). Misclassification in
difference-in-differences models. arXiv:2207.11890.

\bibitem[Freyaldenhoven et al.(2019)]{fhs2019pretrends}
Freyaldenhoven, S., C.~Hansen, and J.~M. Shapiro (2019). Pre-event trends in
the panel event-study design. \emph{American Economic Review} 109(9),
3307--3338.

\bibitem[Gutknecht and Liu(2024)]{gutknecht2026cic}
Gutknecht, D. and C.~Liu (2024). Changes-in-changes for ordered choice models
with underreporting. arXiv:2401.00618.

\bibitem[Hausman et al.(1998)]{hausman1998misclassification}
Hausman, J.~A., J.~Abrevaya, and F.~M. Scott-Morton (1998). Misclassification
of the dependent variable in a discrete-response setting. \emph{Journal of
Econometrics} 87(2), 239--269.

\bibitem[Hotelling(1931)]{hotelling1931}
Hotelling, H. (1931). The generalization of Student's ratio. \emph{Annals of
Mathematical Statistics} 2(3), 360--378.

\bibitem[Kud\^o(1963)]{kudo1963}
Kud\^o, A. (1963). A multivariate analogue of the one-sided test.
\emph{Biometrika} 50(3--4), 403--418.

\bibitem[Lehmann and Romano(2005)]{lehmann2005testing}
Lehmann, E.~L. and J.~P. Romano (2005). \emph{Testing Statistical Hypotheses}
(3rd ed.). New York: Springer.

\bibitem[MacKinnon and Webb(2017)]{mackinnon2017wild}
MacKinnon, J.~G. and M.~D. Webb (2017). Wild bootstrap inference for wildly
different cluster sizes. \emph{Journal of Applied Econometrics} 32(2),
233--254.

\bibitem[Rambachan and Roth(2023)]{rambachan2023credible}
Rambachan, A. and J.~Roth (2023). A more credible approach to parallel trends.
\emph{Review of Economic Studies} 90(5), 2555--2591.

\bibitem[Roth(2022)]{roth2022pretest}
Roth, J. (2022). Pretest with caution: Event-study estimates after testing for
parallel trends. \emph{American Economic Review: Insights} 4(3), 305--322.

\bibitem[Roth et al.(2023)]{roth2023trending}
Roth, J., P.~H.~C. Sant'Anna, A.~Bilinski, and J.~Poe (2023). What's trending
in difference-in-differences? A synthesis of the recent econometrics
literature. \emph{Journal of Econometrics} 235(2), 2218--2244.

\bibitem[Shapiro(1988)]{shapiro1988}
Shapiro, A. (1988). Towards a unified theory of inequality constrained testing
in multivariate analysis. \emph{International Statistical Review} 56(1),
49--62.

\bibitem[Sun and Abraham(2021)]{sunabraham2021}
Sun, L. and S.~Abraham (2021). Estimating dynamic treatment effects in event
studies with heterogeneous treatment effects. \emph{Journal of Econometrics}
225(2), 175--199.

\bibitem[Weisburd et al.(2024)]{weisburd2024crs}
Weisburd, D., D.~B. Wilson, C.~Gill, K.~Kuen, and T.~Zastrow (2024). Are we
underestimating the crime prevention outcomes of community policing? The
importance of crime reporting sensitivity bias. \emph{Journal of Law and
Empirical Analysis} 1(1), 1--17.

\end{thebibliography}

\end{document}
